% !TeX root = ../papers/nc-p3.tex
\section{単純表現の辺を補う自由度と指数成長}
\label{sec:simple-linear-completion-growth}
\begin{lemma}\label{lem:free-tangent-directions}
有限生成代数 $A$ の絶対単純な有限次元表現 $S$ に対し、$m$ 個の非可換変数の自由形式冪級数環
$F=k\langle\!\langle x_1,\ldots,x_m\rangle\!\rangle$ 上の行列表現が $S$ を持ち上げるとする。
その一次方向が $\operatorname{Ext}^1_A(S,S)$ で一次独立なら、$A$ は $m$ 生成自由代数を含む。
\end{lemma}
\begin{proof}
$\mathfrak n=(x_1,\ldots,x_m)$、$D=\dim S$、$E$ を $\mathfrak n^2$ を割った表現の像とする。
Burnsideの定理で定数への射は $E\twoheadrightarrow M_D(k)$ である。核は平方零かつ
$M_D$ 両側部分加群なので $M_D\otimes W$、$W\subset\mathfrak n/\mathfrak n^2$ の形である
（行列単位を左右から掛ける）。$W$ の零化に非零な共変元があれば、その方向への双対数特殊化の
像は定数行列環に同型である。その行列環の持ち上げは恒等に一つの導分を加えたものになり、
行列環の導分は内部なので $S$ の一次変形として自明となる。仮定に反するから $W$ は全体である。
各 $1_D x_j$ を $a_j\in A$ に持ち上げる。$a_{j_1}\cdots a_{j_l}$ の像の最低次は
$1_D x_{j_1}\cdots x_{j_l}$ である。自由語の独立性を関係式の最低次に使うと、$a_j$ の間には
非自明な関係はない。有限生成元による語長で各 $a_j$ は有界なので、$m\ge2$ なら指数成長となる。
\end{proof}

\begin{theorem}\label{thm:simple-linear-completion-growth}
四頂点の二次折り畳みで、空でない互いに隣接しない辺の集合 $J$ を選ぶ（一辺または反対二辺）。
有限次元の非頂点単純 $S$ のうち $J$ を除いた矢作用を $N$ とし、$L_J$ をそれと両立する
$J$ 上の矢作用の全線型空間とする。$e=\dim\operatorname{End}(N)$、$l=\dim L_J$ とする。
折り畳みは $l-e+1$ 生成自由代数を含む。とくに多項式成長なら必ず $l=e$ で、$N$ を固定した
単純な全ての補い方は頂点基底変更と共通次数倍率を除いて一意である。
これは窓の点対応、PI性、Noether性、正則性を仮定しない。
\end{theorem}
\begin{proof}
各二次関係が可変な辺を高々一度通るので、残りを固定すると条件は線型斉次であり、
可変辺を通らない関係は既に成立している。$L_J$ 内の基底変更の接空間は
$\operatorname{End}(N)$ の元 $g$ による $g_{i+1}A_i-A_i g_i$ の像である。
その核は $S$ の自己準同型のスカラーだけで、像の次元は $e-1$ である。
この像の補空間の基底を $D_j$ とすれば $A+\sum_jD_j x_j$ が非可換な自由変数で
全ての関係を正確に満たす。一つの関係項には二変数の積が現れないためである。
一次方向は自己第一Extで独立である。実際頂点冪等元は固定され、内部変形が残りの全ての矢で
零になる条件は $g\in\operatorname{End}(N)$ そのものである。補題を適用する。

$J$ 全体だけを同時に拡大する一次方向は基底変更ではない。ある正長の斉次閉路が $S$ で非零traceを持つ。
実際正次数イデアルの像は全行列環の非零イデアルなので単位を張り、非閉路のtraceは零である。
各閉路は四の倍数長で $J$ を同じ正回数だけ通り、そのtraceの微分が標数零で非零となる。
よって $l-e+1\ge1$、多項式成長なら等号である。$\operatorname{Aut}(N)\times\mathbb G_m$ を
$L_J$ に作用させ、最後の因子は可変辺だけを拡大するものとする。この倍率は代数閉体では
全矢の共通次数倍率と頂点倍率の積で表せる。一周の倍率の積を合わせればよい。
各単純点の軌道の接空間は今や $L_J$ 全体となる。従って各軌道が開であり、線型空間の既約性から
単純点は一軌道である。特にこの結論は補完する辺を選び直す場合にも使える。
\end{proof}

\begin{corollary}\label{cor:two-dimensional-return-pencil}
有限四辺における頂点次元二の安定切断で、戻り辺が存在するとする。その線型解の空間は次元1または2である。
次元2となる必要十分条件は、前後の許す二次曲面が双方とも網であり、その基底の滑らかな空間三次曲線が
一致することである。この場合、二次元の空間 $H$ を用い、戻り文字を $v\in\operatorname{Sym}^3 H^*$ と
同定すると、全解は $a\in H$ を動かした偏極 $X_a(v)=v(a,-,-)$ となる。全ての非零解の行列式は
階数3の錐であり、この窓は指数成長なので正則候補から除ける。
それ以外では唯一の共通二次曲面と両向きの行列表示が解を定める。
\end{corollary}
\begin{proof}
$K_+$ と $K_-$ を戻り文字 $V$ での前からの行解と後ろからの列解の空間とする。
命題\ref{prop:two-dimensional-quadric-successors}からそれぞれ次元2または3で、行列階数1の元を持たない。
$L$ を戻り解全体とすると全非零解は単純を与える。よって全ての非零な出力共変元 $u$ と $X\ne0$ で
$uX\ne0$、任意の非零な入力 $w$ で $Xw\ne0$ である。
$u\otimes X\mapsto uX$ の核は可分解元を含まないので、Segreの次元から $\dim K_+\ge\dim L+1$。
列版も同じで、$\dim L\le2$ である。次元2なら両方3で、両縮約の核の直線は二行二列可逆行列で張られる。
よって三つの二次元空間（出力双対、$L$、入力）の基底を選んで両直線をそれぞれ交代テンソルとできる。
戻り解の普遍四テンソルはこの三因子の隣接互換二つについて対称になり、全対称三乗を通る。
$V\to\operatorname{Sym}^3 H^*$ は単射である。実際三つの行をまとめた写像の単射性は
命題\ref{prop:two-dimensional-quadric-successors}で示した。次元から同型である。
各 $a$ に対し三次式の偏極は全二次式へ全射し、対称二行二列の行列式は非退化な三変数二次式である。
三次曲線は双方で純三乗の軌跡となる。

逆に同じ三次曲線 $C$ を持つ二網を考える。各側で $[w]\mapsto\ker(B(-)w)$ は
$\mathbb P^1\simeq C\subset\mathbb P(V)$ の完全三次表示であった。
従って各側の三行は、このパラメータ表示の全ての一次シジジーを与える。
具体的に $[1:t:t^2:t^3]$ の斉次化の四つの三次単項式に対する一次シジジーは
隣り合う項の三つの差で張られ、次元3である。左右のパラメータ表示は射影一次変換で揃う。
これで $V$ を二元三次式の空間、両側の行・列を偏極の三行・三列とする同じ表示を取れる。
二つの独立な $a$ による偏極はいずれも両制約を満たし、先の上界からこれで全解である。
安定切断の自己準同型はスカラーだけなので、定理\ref{thm:simple-linear-completion-growth}より二解なら自由二生成となる。
一網でない側があればその側の二次曲面は既に一つである。異なる二曲線を持つ二網でも、
独立な二共通二次曲面はない。各曲面は階数3以上で既約、二曲面の交叉の次数4は二曲線の次数和6に足りない。
可逆な接続と向きの検査は定理\ref{thm:finite-two-dimensional-quadrics}で与えた。
\end{proof}
